We introduce QuaRot: a method which uses Hadamard matrices to eliminate outliers in the activations and KV cache of pre-trained LLMs, enabling end-to-end 4-bit quantization for the first time (to the best of our knowledge). Quantizing \llamaBig to 4 bits with QuaRot maintains 99\% of the downstream task performance of the FP16 baseline, with a 2.16$\times$ speedup on RTX 3090 GPUs during the prefill stage (and up to 3.39$\times$ memory saving during the decoding stage). Quantizing all \llama models to 6 and 8 bits is lossless. 

Opportunities to build on QuaRot include quantizing the residuals and extending the method to mixture-of-experts architectures. In terms of hardware, end-to-end INT4 inference with QuaRot could be exploited to give similar speedups as that of the recently announced NVIDIA B200 GPU architecture, while being much cheaper to implement compared to the floating point (FP4) format.
